%!TEX program = xelatex
% !TEX encoding = ISO-8859-1
\def\ano{2020}
\def\semestre{1}
\def\disciplina{\'Algebra 1}
\def\turma{C}
\def\numerosemana{07}

\documentclass[12pt]{exam}

\usepackage{caption}
\usepackage{amssymb}
\usepackage{amsmath,amsfonts,amsthm,amstext}
\usepackage[brazil]{babel}
% \usepackage[latin1]{inputenc}
\usepackage{graphicx}
\graphicspath{{/home/jfreitas/GitHub_Repos/Algebra-1/Pictures/}{D:/OneDrive - unb.br/imagens-latex/}}
\usepackage{enumitem}
\usepackage{multicol}
\usepackage{answers}
\usepackage{tikz,ifthen}
\usetikzlibrary{lindenmayersystems}
\usetikzlibrary[shadings]
\Newassociation{solucao}{Solution}{ans}
\newtheorem{exercicio}{}

\setlength{\topmargin}{-1.0in}
\setlength{\oddsidemargin}{0in}
\setlength{\textheight}{10.1in}
\setlength{\textwidth}{6.5in}
\setlength{\baselineskip}{12mm}

\extraheadheight{0.7in}
\firstpageheadrule
\runningheadrule
\lhead{
        \begin{minipage}[c]{1.7cm}
        \includegraphics[width=1.7cm]{unb.pdf}
        \end{minipage}%
        \hspace{0pt}
        \begin{minipage}[c]{4in}
          {Universidade de Bras{\'\i}lia} --
          {Departamento de Matem{\'a}tica}
\end{minipage}
\vspace*{-0.8cm}
}
% \chead{Universidade de Bras{\'\i}lia - Departamento de Matem{\'a}tica}
% \rhead{}
% \vspace*{-2cm}

\extrafootheight{.5in}
\footrule
\lfoot{\disciplina\ - \semestre$^o$/\ano}
\cfoot{}
\rfoot{P\'agina \thepage\ de \numpages}

\newcounter{exercicios}
\renewcommand{\theexercicios}{\arabic{exercicios}}

\newenvironment{questao}[1]{
\refstepcounter{exercicios}
\ifx&#1&
\else
   \label{#1}
\fi
\noindent\textbf{Exerc{\'\i}cio {\theexercicios}:}
}

\newcommand{\resp}[1]{
\noindent{\bf Exerc{\'\i}cio #1: }}
\newcommand{\nsub}{\varsubsetneq}
\newcommand{\vaz}{\emptyset}
\newcommand{\im}{{\rm Im\,}}
\newcommand{\sub}{\subseteq}
\newcommand{\n}{\mathbb{N}}
\newcommand{\z}{\mathbb{Z}}
\newcommand{\rac}{\mathbb{Q}}
\newcommand{\real}{\mathbb{R}}
\newcommand{\complex}{\mathbb{C}}
\newcommand{\cp}[1]{\mathbb{#1}}
\newcommand{\ch}{\mbox{\textrm{car\,}}\nobreak}
\newcommand{\vesp}[1]{\vspace{ #1  cm}}
\newcommand{\compcent}[1]{\vcenter{\hbox{$#1\circ$}}}
\newcommand{\comp}{\mathbin{\mathchoice
{\compcent\scriptstyle}{\compcent\scriptstyle}
{\compcent\scriptscriptstyle}{\compcent\scriptscriptstyle}}}

\begin{document}
    \begin{center}
    {\Large\bf \disciplina\ - Turma \turma\ -- \semestre$^{o}$/\ano} \\ \vspace{9pt} {\large\bf
        Lista de Exerc{\'\i}cios -- Semana \numerosemana}\\ \vspace{9pt} Prof. Jos{\'e} Ant{\^o}nio O. Freitas
    \end{center}
    \hrule

    \vspace{.6cm}

    \begin{center}
    \textit{Nota\c{c}\~oes:}
\end{center}
\begin{multicols}{2}
    \begin{enumerate}[label={\roman*})]
        \item $\n^* = \n - \{0\}$
        \item $\real^*_+ = \{x \in \real \mid x > 0\}$
        \item $\real^*_- = \{x \in \real \mid x < 0\}$
        \item $\real_+ = \{x \in \real \mid x \ge 0\}$
        \item $\real_- = \{x \in \real \mid x \le 0\}$
        \item $]c,d[\ = (c,d) = \{x \in \real \mid c < x < d\}$
        \item $[c,d] = \{x \in \real \mid c \le x \le d\}$
        \item $[c,d[\ = [c, d) = \{x \in \real \mid c \le x < d\}$
        \item $]c,d] = (c, d] = \{x \in \real \mid c < x \le d\}$
        \item Em $\z_m$ vamos denotar a opera\c{c}\~ao de soma $\oplus$ por $+$.
        \item Em $\z_m$ vamos denotar a opera\c{c}\~ao de multiplica\c{c}\~ao $\otimes$ por $\cdot$.
    \end{enumerate}
\end{multicols}

\vspace{.6cm}

\questao{} Considere a fun\c{c}\~ao $f : \real \to \real$ dada por $f(x) = \mid x - 2\mid$. Encontre:
\begin{multicols}{3}
    \begin{enumerate}[label={\alph*})]
        \item $f(\{1\})$
        \item $f(\{-\sqrt{2}, 3\})$
        \item $f([-2,2])$
        \item $f((-3,5))$
        \item $f^{-1}(\{3\})$
        \item $f^{-1}(\{-3,5\})$
        \item $f^{-1}([0,2])$
        \item $f^{-1}([-3,3])$
        \item $f^{-1}(\real^*_-)$
    \end{enumerate}
\end{multicols}

\vspace{.3cm}

\questao{} Seja $g : \real \to \real$ dada por
\[
    g(x) = \begin{cases}
        x^2,& \mbox{ se } x \le 0\\
        \sqrt[3]{x}, & \mbox{ se } x > 0.
    \end{cases}
\]
Encontre:
\begin{multicols}{2}
    \begin{enumerate}[label={\alph*})]
        \item $g([-1,8])$
        \item $g(\real_+)$
        \item $g^{-1}([-1,16])$
        \item $g(\real_-)$
        \item $g^{-1}([1,25])$
        \item $g^{-1}(\real^*_-)$
    \end{enumerate} 
\end{multicols}


\vspace{.3cm}

\questao{} Seja $f(x) = x^4$ e $g(x) = x^7$. Verifique que $(f\circ g)(x) = (g\circ f)(x)$.

\vspace{.3cm}

\questao{} Dadas as fun\c{c}\~oes $f(x) = 3x + m$ e $g(x) = ax + 2$, determine condi\c{c}\~oes sobre $a$ e $m$ para que $(f\circ g)(x) = (g\circ f)(x)$.

\newpage

\questao{} Sejam $f$, $g$ e $h$ fun\c{c}\~oes de $\real$ em $\real$ definidas por $f(x) = x + 2$, $g(x) = x^2 - 1$ e $h(x) = 3x$, respectivamente.
\begin{enumerate}[label={\alph*})]
    \item Determine $f \circ g$, $f \circ h$, $g \circ f$, $g \circ h$ e $h \circ g$.
    \item Verifique que $(f \circ g)\circ h = f \circ (g \circ h)$.
\end{enumerate}

\vspace{.3cm}

\questao{} Dada as fun\c{c}\~oes
\[
    f(x) = \begin{cases}
        1, & \mbox{ se } x < 0\\
        2x^2, & \mbox{ se } 0 \le x \le 1\\
        0, & \mbox{ se } x > 1
    \end{cases} \qquad g(x) = \begin{cases}
        x, & \mbox{ se } x < 0\\
        0, & \mbox{ se } 0 \le x \le 1\\
        1, & \mbox{ se } x > 1.
    \end{cases}
\]
Determine $f\circ g$.

\vspace{.3cm}

\questao{} Dada as fun\c{c}\~oes
\[
    f(x) = \begin{cases}
        x^2 + 2, & \mbox{ se } x \le -1\\
        \dfrac{1}{x - 2}, & \mbox{ se } -1 < x < 1\\
        4 - x^2, & \mbox{ se } x \ge 1
    \end{cases} \qquad g(x) = 2 - 3x.
\]
Determine $f\circ g$ e $g \circ f$.

\vspace{.3cm}

\questao{} Seja $f: \real^2 \to \real$ dada por $f(x,y) = xy$.
\begin{enumerate}[label={\alph*})]
    \item $f$ {\'e} injetora?
    \item $f$ {\'e} sobrejetora?
    \item Obter $f^{-1}({0})$.
    \item Obter $f([0,1]\times [0,1])$.
\end{enumerate}

\vspace{.3cm}

\questao{} Seja $f : A \to [-9,-1)$ dada por $f(x) = \dfrac{4x + 3}{3 - x}$.
\begin{enumerate}[label={\alph*})]
    \item Determine $A$.
    \item Mostre que $f$ \'e injetora.
    \item \'E verdade que $f$ \'e sobrejetora?
\end{enumerate}

\vspace{.3cm}

\questao{} Seja $f : A \to (1,10]$ dada por $f(x) = \dfrac{4 - 11x}{4 - 2x}$.
\begin{enumerate}[label={\alph*})]
    \item Determine $A$.
    \item Mostre que $f$ \'e injetora.
    \item \'E verdade que $f$ \'e sobrejetora?
\end{enumerate}

\vspace{.3cm}

\questao{} Considere a fun{\c c}{\~a}o $f : \z \times \z \to \z \times \z$ tal que $f(x,y) = (2x + 3, 4y + 5)$. Prove que $f$ {\'e} injetora. Verifique se $f$ {\'e} bijetora.

\vspace{.3cm}

\questao{} Mostre que a fun\c{c}\~ao $f : \z \to \z$ dada por $f(n) = 2n$, \'e injetora mas n\~ao \'e sobrejetora.

\vspace{.3cm}

\questao{} Mostre que a fun\c{c}\~ao $h : \real^2 \to \real^2$ tal que $h(x, y) = (\sqrt[3]{x}, y^5)$ \'e sobrejetora.
\vspace{.3cm}

\questao{} Considere a fun{\c c}{\~a}o $g : \z_5 \times \z_9 \to \z_5 \times \z_9$ tal que $f(\overline{x},\overline{y}) = (\overline{2} \overline{x} + \overline{3}, \overline{4}\overline{y} + \overline{5})$.
\begin{enumerate}[label={\alph*})]
    \item $g$ \'e injetora?
    \item $g$ \'e sobrejetora?
\end{enumerate}

\vspace{.3cm}

\questao{} Considere a fun{\c c}{\~a}o $g : \z_5 \times \z_7 \to \z_5 \times \z_7$ tal que $g(\overline{x},\overline{y}) = (\overline{x} + \overline{3}, \overline{6}\overline{y} + \overline{2})$. Determine a fun\c{c}\~ao inversa de $g$.

\vspace{.3cm}

\questao{} Mostre que a fun{\c c}{\~a}o $f : \real \to \real$ definida por $f(x) = ax + b$, com $a$ e $b$ constantes reais, $a \ne 0$, {\'e} uma bije{\c c}{\~a}o. Obter $f^{-1}$.

\vspace{.3cm}

\questao{} Mostrar que $f : \real - \left\{-\dfrac{d}{c}\right\} \to \real  - \left\{\dfrac{a}{c}\right\}$ dada por $f(x) =  \dfrac{ax + b}{cx + d}$, onde $a$, $b$, $c$, $d$ s{\~a}o n{\'u}meros reais constantes, $ad - bc \ne 0$, {\'e} uma bije{\c c}{\~a}o. Descrever a fun{\c c}{\~a}o $f^{-1}$.

\vspace{.3cm}

\questao{} Achar uma fun{\c c}{\~a}o $f : A \to B$, com $A$ e $B$ subconjuntos de $\real$, para cada caso abaixo:
\begin{enumerate}[label={\alph*})]
    \item $A = \real$, $B \nsub \real$ e $f$ injetora e n{\~a}o sobrejetora.
    \item $A \nsub \real$, $B = \real$ e $f$ injetora e n{\~a}o sobrejetora.
    \item $A = \real$, $B \nsub \real$ e $f$ sobrejetora e n{\~a}o injetora.
    \item $A \nsub \real$, $B = \real$ e $f$ sobrejetora e n{\~a}o injetora.
\end{enumerate}

\vspace{.3cm}

\questao{} Dois conjuntos, $A$ e $B$, s\~ao \textbf{equipotentes} quando $A = B = \emptyset$ ou existe $f : A \to B$ fun\c{c}\~ao   bijetora. Mostre, em cada caso, que os conjuntos $A$ e $B$ s\~ao equipotentes, encontrado uma fun\c{c}\~ao bijetora $f : A \to B$ em cada caso:
\begin{enumerate}[label={\alph*})]
    \item $A = \n$ e $B = \n^*$
    \item $A = \z$ e $B = \n$
    \item $A = \real$ e $B = \real$
    \item $A = ]-1, 1[$ e $B = ]c,d[$ com $c$ e $d$ n\'umeros reais e $c < d$.
\end{enumerate}

\newpage

\questao{} Classificar (se poss{\'\i}vel) em injetora ou sobrejetora as seguintes fun{\c c}{\~o}es de $\real$ em $\real$.

\begin{multicols}{2}
    \begin{enumerate}[label={\alph*})]
        \item $f(x) = x^3$
        \item $f(x) = x^2 - 5x - 6$
        \item $f(x) = 2^x$
        \item $f(x) = | \sin x |$
        \item $f(x) = x + | x |$
        \item $f(x) = x + 3$
        \item $f(x) = \mid x - 1\mid$
        \item $f(x) = \dfrac{1}{x}$
        \item $f(x) = 1 - x^2$
        \item $f(x) = |x|(x - 1)$
    \end{enumerate}
\end{multicols}

\vspace{.3cm}

\questao{} Seja $f : A \to B$ e $g : B \to A$ fun\c{c}\~oes tais que $g \circ f = i_A$. Quais das afirma\c{c}\~oes seguintes s\~ao verdadeiras?
\begin{enumerate}[label={\alph*})]
    \item $g = f^{-1}$
    \item $f$ \'e sobrejetora
    \item $f$ \'e injetora
    \item $g$ \'e sobrejetora
    \item $g$ \'e injetora
\end{enumerate}


\vspace{.3cm}

\questao{} Sejam $f : A \to B$, $g : A \to B$ e $h : B \to C$ fun\c{c}\~oes. Prove que se $h$ \'e injetora e $h \circ g = h \circ f$, ent\~ao $g = f$.

\vspace{.3cm}

\questao{} Seja $f : A \to B$ uma fun\c{c}\~ao e sejam $P$ e $Q$ subconjuntos de $A$. Mostre que:
\begin{enumerate}[label={\alph*})]
    \item Se $P\sub Q$, ent{\~a}o $f(P)\sub f(Q)$.
    \item $f(P\cup Q) = f(P)\cup f(Q)$.
    \item $f(P\cap Q)\sub f(P)\cap f(Q)$.
    \item Se $f$ {\'e} injetora, ent{\~a}o $f(P\cap B) =  f(P)\cap f(Q)$.
    \item $f$ {\'e} bijetora se, e somente se, $f(P^C) = (f(P))^C$ para todo $P \sub A$. \textit{(Aqui $P^C$ \'e o complementar de $P$ em rela\c{c}\~ao \`a $A$.)}
\end{enumerate}

\vspace{.3cm}

\questao{} Seja $f : A \to B$ uma fun{\c c}{\~a}o e sejam $P \sub
A$ e $X, Y\sub B$. Mostre que:
\begin{enumerate}[label={\alph*})]
    \item Se $X\sub Y$, ent{\~a}o $f^{-1}(X)\sub f^{-1}(Y)$.
    \item $f^{-1}(X\cup Y)=f^{-1}(X)\cup f^{-1}(Y)$.
    \item $f^{-1}(X\cap Y)= f^{-1}(X)\cap f^{-1}(Y)$.
    \item $P\sub f^{-1}(f(P))$.
    \item $f(f^{-1}(X))= X \cap \mbox{Im}f$ e conclua que se $f$ {\'e} sobrejetora ent{\~a}o
    $f(f^{-1}(X))=X$.
    \item $f^{-1}(X^C) = (f^{-1}(X))^C$.
    \item $f$ \'e sobrejetora se, e somente se, $f^{-1}(T) \ne \emptyset$ para todo $T \sub B$.
\end{enumerate}

\vspace{.3cm}

\questao{} Se as fun{\c c}{\~o}es $f : A \to B$ e $g : B\to A$ s{\~a}o
tais que $g\circ f$ {\'e} injetora, ent{\~a}o $f$ {\'e} injetora.

\vspace{.3cm}

\questao{} Se as fun{\c c}{\~o}es $f : A \to B$ e $g : B\to A$ s{\~a}o
tais que $g\circ f$ {\'e} sobrejetora, ent{\~a}o g {\'e} sobrejetora.

\vspace{.3cm}

\questao{} Mostrar que toda fun{\c c}{\~a}o injetora (sobrejetora) de um conjunto finito em si mesmo {\'e} tamb{\'e}m sobrejetora (injetora).

\end{document}